\lysh{31089}{2}
\begin{alist}
\item  Το πεδίο ορισμού της συνάρτησης $f$ είναι το σύνολο που περιέχει τις τιμές για τις οποίες ορίζεται το κλάσμα, δηλαδή για $x \neq 0$. Συνεπώς, $D_f = \mathbb{R} - \{0\}$.
\item Βρίσκουμε την πρώτη παράγωγο της συνάρτησης.
Παραγωγίζουμε τη συνάρτηση ως άθροισμα δύο συναρτήσεων:
\[ f'(x) = \left(4 \cdot \frac{1}{x}\right)' - (1)' = 4 \cdot \left(-\frac{1}{x^2}\right) - 0 = 4 \cdot \left(-\frac{1}{x^2}\right) = -\frac{4}{x^2}. \]
Άρα, $f'(x) = -\frac{4}{x^2}$ για κάθε $x \in \mathbb{R} - \{0\}$.
\item Από το β) ερώτημα έχουμε $f'(x) = -\frac{4}{x^2}$.
Άρα, για $x=-2$ έχουμε:
\[ f'(-2) = -\frac{4}{(-2)^2} = -\frac{4}{4} = -1. \]
\end{alist}
\vspace{6mm}
\noindent\rule{\textwidth}{0.4pt}
\vspace{6mm}

